\documentclass[11pt]{article}
\usepackage[margin=1in]{geometry}
\usepackage{amsmath,amssymb,amsthm}

\newtheorem{definition}{Definition}
\newtheorem{proposition}{Proposition}
\newtheorem{remark}{Remark}

\newcommand{\phig}{\varphi}
\newcommand{\mustar}{\mu_\star}
\newcommand{\Ecoh}{E_{\mathrm{coh}}}
\newcommand{\Epass}{E_{\mathrm{pass}}}
\newcommand{\Bpow}{B_{\mathrm{pow}}}

\title{Charged Lepton Mass Ratios from the Combinatorics of the 3-Cube\\
\large Integer Backbone and Canonical Refinement Structure}

\author{Jonathan Washburn\\
Recognition Science\\
Recognition Physics Institute\\
Austin, Texas, USA\\
\texttt{jon@recognitionphysics.org}}

\date{March 2026}

\begin{document}
\maketitle

\begin{abstract}
This paper isolates the charged-lepton ratio sector of the Recognition Science mass framework and states its claim boundary in publishable form. The backbone is purely combinatorial. From the three-cube \(Q_3\) one obtains the fixed integer vocabulary
\[
V=8,\qquad E=12,\qquad F=6,\qquad A=1,\qquad \Epass=11,\qquad W=17.
\]
For the charged leptons, the common charge coordinate is
\[
Z_\ell=(6Q)^2+(6Q)^4=1332\qquad (Q=-1),
\]
so the charge band cancels in family ratios at the common anchor scale. With the lepton rungs
\[
r_e=2,\qquad r_\mu=13,\qquad r_\tau=19,
\]
the integer backbone gives
\[
\frac{m_\mu}{m_e}=\phig^{11},\qquad \frac{m_\tau}{m_\mu}=\phig^{6},
\]
where \(\phig=(1+\sqrt5)/2\). To compare against reference masses at precision level, one adds the canonical refinement structure
\[
S_{e\to\mu}=\Epass+\frac{1}{4\pi}-\alpha^2,
\qquad
S_{\mu\to\tau}=F-\frac{2W+D}{2}\,\alpha,
\qquad D=3,
\]
using the framework value \(\alpha^{-1}\approx137.0348851720\). This yields
\[
\frac{m_\mu}{m_e}\approx206.76814,
\qquad
\frac{m_\tau}{m_\mu}\approx16.81559,
\]
within about \(7\times10^{-7}\) and \(9\times10^{-5}\) relative error of standard reference ratios. The leading integer backbone is treated here as derived from the \(Q_3\) combinatorics. The sub-leading correction terms are not claimed as closed first-principles derivations in this paper; they are stated honestly as the current canonical refinement structure of the lepton chain.
\end{abstract}

\section{Introduction}

The Standard Model accommodates the charged lepton masses through Yukawa couplings but does not derive the hierarchy
\[
m_e \ll m_\mu \ll m_\tau
\]
from a smaller discrete substrate. The Recognition Science program proposes that the hierarchy is controlled by a golden-ratio ladder and a small integer vocabulary supplied by the three-cube \(Q_3\). The present paper is deliberately narrower than a full mass-spectrum paper. Its purpose is to isolate the charged-lepton ratios, separate the integer backbone from the small correction terms, and state exactly what is and is not being claimed.

That separation matters. The current framework already supports a clean integer story for the charged leptons: the family shares a common charge coordinate, the lepton sector shares a common yardstick, and the family rungs are fixed by the generation-torsion pattern. Once those facts are written at the common anchor scale, the mass ratios reduce to powers of \(\phig\). Precision comparison then introduces a second layer: small sub-leading corrections to the two family steps. Those corrections are numerically successful, but their full first-principles derivation is not claimed as closed here.

The paper therefore makes one strong claim and one weaker claim. The strong claim is that the charged-lepton ratio backbone is a direct consequence of the \(Q_3\) combinatorics. The weaker claim is that the current canonical refinement formulas push the backbone close to the observed ratios without species-by-species fitting.

\section{Status words and scope}

Three status words are used throughout.

\paragraph{Derived backbone.} A statement belongs to the derived backbone when it follows from the fixed cube vocabulary, the lepton charge map, and the rung construction at the common anchor scale.

\paragraph{Canonical refinement.} A statement belongs to the canonical refinement layer when it is part of the current lepton-chain formula used for numerical comparison, but this paper does not claim that it has already been reduced to a closed first-principles theorem.

\paragraph{Reference display.} A statement belongs to the reference display layer when it uses standard charged-lepton mass values only for comparison with the framework output.

This paper does not discuss quarks, neutrinos, or the full SI seam. It uses only the minimal anchor-scale machinery needed to turn the lepton family into a ratio problem. The fuller anchor-scale law, sector yardsticks, and calibration protocol are given in the companion backbone paper \emph{The Recognition Science Fermion Mass Backbone at the Anchor Scale}.

\section{Minimal methods summary from the backbone paper}

At the common charged-fermion anchor \(\mustar\), the structural mass law is
\[
  m_f^{\mathrm{RS}}(\mustar)
  = A_{s(f)}\,\phig^{\,r_f - 8 + \mathrm{gap}(Z_f)}.
\]
Here \(s(f)\) is the sector, \(A_s\) is the sector yardstick, \(r_f\) is the rung, and \(Z_f\) is the integerized charge coordinate. For the lepton sector,
\[
A_{\mathrm{L}} = 2^{\Bpow(\mathrm{L})}\Ecoh\,\phig^{r_0(\mathrm{L})}
=2^{-22}\Ecoh\,\phig^{62},
\qquad
\Ecoh=\phig^{-5},
\]
and the common anchor is \(\mustar=182.201\,\mathrm{GeV}\).

The only feature of this law needed in the present paper is the following cancellation principle.

\begin{proposition}[Equal-\(Z\) ratio cancellation]
If two fermions lie in the same sector and share the same charge coordinate \(Z\), then at the common anchor scale
\[
\frac{m_f^{\mathrm{RS}}(\mustar)}{m_g^{\mathrm{RS}}(\mustar)}=\phig^{r_f-r_g}.
\]
\end{proposition}

\begin{proof}
The common sector factor \(A_s\), the octave offset \(-8\), and the common charge band \(\mathrm{gap}(Z)\) cancel in the ratio. Only the rung difference remains.
\end{proof}

The charged leptons satisfy exactly this equal-\(Z\) condition. Using the lepton charge map
\[
Z=(6Q)^2+(6Q)^4,
\]
and the charged-lepton electric charge \(Q=-1\), one obtains
\[
Z_\ell=(-6)^2+(-6)^4=36+1296=1332.
\]
Thus the charged-lepton ratios reduce to rung differences at \(\mustar\).

\section{The \(Q_3\) integer backbone}

The three-cube supplies the entire small-integer vocabulary used in the lepton ratio sector:
\[
V=8,
\qquad
E=12,
\qquad
F=6,
\qquad
A=1,
\qquad
\Epass=E-A=11,
\qquad
W=\Epass+F=17.
\]
The quantities \(\Epass\) and \(W\) are the relevant generation-torsion integers for the charged-lepton family.

\begin{proposition}[Lepton rungs]
The charged-lepton rungs are
\[
r_e=2,
\qquad
r_\mu=2+\Epass=13,
\qquad
r_\tau=2+W=19.
\]
Hence the two consecutive family steps are
\[
r_\mu-r_e=\Epass=11,
\qquad
r_\tau-r_\mu=W-\Epass=F=6.
\]
\end{proposition}

\begin{proof}
The lepton baseline is \(r_e=2\). The second-generation torsion adds \(\Epass=11\), giving \(r_\mu=13\). The third-generation torsion adds \(W=17\), giving \(r_\tau=19\). The consecutive differences are therefore \(13-2=11\) and \(19-13=6\). Since \(W=\Epass+F\), the second difference is \(W-\Epass=F\).
\end{proof}

Combining the rung proposition with equal-\(Z\) cancellation gives the integer ratio backbone.

\begin{proposition}[Integer backbone ratios]
At the common anchor scale,
\[
\frac{m_\mu^{\mathrm{RS}}(\mustar)}{m_e^{\mathrm{RS}}(\mustar)}=\phig^{11},
\qquad
\frac{m_\tau^{\mathrm{RS}}(\mustar)}{m_\mu^{\mathrm{RS}}(\mustar)}=\phig^{6}.
\]
Numerically,
\[
\phig^{11}\approx 199.0050,
\qquad
\phig^6\approx 17.9443.
\]
\end{proposition}

The integer backbone already puts the family hierarchy in the correct qualitative order, but it is not yet a precision description. The next section states the canonical refinement layer used for numerical comparison.

\section{Canonical refinement structure}

The lepton chain uses two refined step exponents,
\[
S_{e\to\mu}=\Epass+\varepsilon_{e\mu},
\qquad
S_{\mu\to\tau}=F+\varepsilon_{\mu\tau},
\]
with the current canonical refinement terms
\begin{align*}
\varepsilon_{e\mu} &= \frac{1}{4\pi}-\alpha^2, \\
\varepsilon_{\mu\tau} &= -\frac{2W+D}{2}\,\alpha,
\qquad D=3.
\end{align*}
Equivalently,
\begin{align}
S_{e\to\mu} &= \Epass+\frac{1}{4\pi}-\alpha^2,
\label{eq:step_emu} \\
S_{\mu\to\tau} &= F-\frac{2W+D}{2}\,\alpha.
\label{eq:step_mutau}
\end{align}
Because \(2W+D=2\cdot17+3=37\), the second step is also
\[
S_{\mu\to\tau}=6-\frac{37}{2}\,\alpha.
\]

\begin{remark}[Claim boundary for the small corrections]
The role of equations \eqref{eq:step_emu} and \eqref{eq:step_mutau} in this paper is explicit and limited. Their leading integers, \(\Epass=11\) and \(F=6\), belong to the derived backbone. The small terms \(1/(4\pi)\), \(-\alpha^2\), and \(-(2W+D)\alpha/2\) are used here as the current canonical refinement structure of the lepton chain. This paper does not present them as closed first-principles theorems. That boundary is deliberate. It keeps the precision numerics visible without overstating the derivational status.
\end{remark}

A useful interpretation remains available even under that stricter claim boundary. The first family step is dominated by the passive-edge count \(\Epass=11\), while the second is dominated by the face count \(F=6\). The charged-lepton hierarchy is therefore controlled first by the one-dimensional edge network and then by the two-dimensional face structure of the cube. The refinement terms are small perturbations around that integer skeleton.

\section{Numerical evaluation of the ratios}

For numerical evaluation, this paper uses the framework value
\[
\alpha^{-1}\approx137.0348851720,
\qquad
\alpha\approx0.007297411887.
\]
Substituting this value into the canonical step formulas gives
\begin{align*}
S_{e\to\mu} &\approx 11+\frac{1}{4\pi}-\alpha^2
\approx 11.079524219, \\
S_{\mu\to\tau} &\approx 6-\frac{37}{2}\alpha
\approx 5.864997880.
\end{align*}
The resulting refined ratios are
\begin{align}
\frac{m_\mu}{m_e} &\approx \phig^{11.079524219}
\approx 206.76814,
\label{eq:ratio_emu_num} \\
\frac{m_\tau}{m_\mu} &\approx \phig^{5.864997880}
\approx 16.81559.
\label{eq:ratio_mutau_num}
\end{align}

For comparison, using the standard charged-lepton reference masses
\[
m_e=0.51099895\,\mathrm{MeV},
\qquad
m_\mu=105.6583755\,\mathrm{MeV},
\qquad
m_\tau=1776.86\,\mathrm{MeV},
\]
one obtains the reference ratios
\[
\frac{m_\mu}{m_e}\approx206.76828,
\qquad
\frac{m_\tau}{m_\mu}\approx16.81703.
\]

\begin{center}
\begin{tabular}{|l|c|c|c|c|}
\hline
Step & Integer backbone & Canonical step & Predicted ratio & Reference ratio \\
\hline
\(e\to\mu\) & \(11\) & \(11.079524219\) & \(206.76814\) & \(206.76828\) \\
\hline
\(\mu\to\tau\) & \(6\) & \(5.864997880\) & \(16.81559\) & \(16.81703\) \\
\hline
\end{tabular}
\end{center}

The corresponding relative discrepancies are about
\[
-6.9\times10^{-7}
\qquad\text{and}\qquad
-8.6\times10^{-5}.
\]
For the purposes of this paper, the important point is not the last displayed digit. It is the separation of roles: the large hierarchy is carried by the integer backbone, while the small correction terms move the ratios from rough order-of-magnitude placement to near-reference numerical agreement.

\section{What is and is not being claimed}

This paper claims the following.

First, the charged-lepton family admits an equal-\(Z\) reduction at the common anchor scale, so family ratios are rung ratios. Second, the rung data give an unambiguous integer backbone:
\[
11\quad\text{for }e\to\mu,
\qquad
6\quad\text{for }\mu\to\tau.
\]
Third, the current canonical refinement structure produces numerically accurate ratios without species-by-species mass fitting.

This paper does not claim the following.

It does not claim a closed first-principles derivation of the small refinement terms in equations \eqref{eq:step_emu} and \eqref{eq:step_mutau}. It does not claim a quark-sector derivation, a neutrino derivation, or a full SI calibration analysis. It does not claim that every digit in the refined ratios is already locked by theorem. The point of the present paper is more modest and cleaner: the integer backbone is clear, the refinement structure is explicit, and the open derivational frontier is not hidden.

That distinction is useful scientifically. If a later derivation confirms the present refinement terms, then this paper becomes the clean minimal lepton-ratio derivation. If instead the refinement terms change, the integer cube backbone still survives as the robust core statement and only the precision layer needs revision.

\appendix

\section{Numerical reproduction recipe}

The main numerical outputs of the paper can be reproduced from the following sequence.

Start with
\[
\phig=\frac{1+\sqrt5}{2},
\qquad
D=3,
\qquad
\Epass=11,
\qquad
F=6,
\qquad
W=17,
\qquad
\alpha^{-1}=137.0348851720.
\]
Then compute
\[
\alpha=\frac{1}{137.0348851720}\approx0.007297411887.
\]
Next evaluate
\[
S_{e\to\mu}=11+\frac{1}{4\pi}-\alpha^2\approx11.079524219,
\]
and
\[
S_{\mu\to\tau}=6-\frac{37}{2}\alpha\approx5.864997880.
\]
Finally compute
\[
\phig^{S_{e\to\mu}}\approx206.76814,
\qquad
\phig^{S_{\mu\to\tau}}\approx16.81559.
\]
No charged-lepton mass value is used on the right-hand side of these two predictions.

\section{Formula provenance inside the framework}

\begin{center}
\begin{tabular}{|l|l|l|}
\hline
Formula or datum & Role in this paper & Status here \\
\hline
\(m_f^{\mathrm{RS}}(\mustar)=A_s\phig^{r_f-8+\mathrm{gap}(Z_f)}\) & anchor-scale mass law & imported backbone law \\
\hline
\(Z_\ell=1332\) & equal-\(Z\) cancellation for leptons & derived backbone \\
\hline
\(r_e=2,\ r_\mu=13,\ r_\tau=19\) & family rung data & derived backbone \\
\hline
\(m_\mu/m_e=\phig^{11},\ m_\tau/m_\mu=\phig^6\) & integer family ratios & derived backbone \\
\hline
\(S_{e\to\mu}=11+1/(4\pi)-\alpha^2\) & refined first step & canonical refinement \\
\hline
\(S_{\mu\to\tau}=6-(37/2)\alpha\) & refined second step & canonical refinement \\
\hline
\(\alpha^{-1}\approx137.0348851720\) & global dimensionless input & imported constant-layer value \\
\hline
Reference masses in MeV & external comparison only & reference display \\
\hline
\end{tabular}
\end{center}

\end{document}
